% !TEX root = ../PhD Thesis.tex
\chapter{Mitochondria, SGs and Inflammation: A novel regulatory triade}
\label{chap:ChapterVI}

Since senescent cells exhibit cytoplasmic accumulation of mitochondrial dsRNA that drives inflammation63, and SGs have been reported to counteract senescence14, we propose a working model (Figure 8D). In the absence of SGs, stress promotes leakage of mitochondrial RNA and DNA into the cytoplasm, where they activate cGAS–STING and RIG-I/MDA5/MAVS pathways63. These innate immune pathways normally detect viral nucleic acids (like dsRNA) and trigger interferon and inflammatory responses, ultimately promoting SASP64. Despite its cellular intrinsic origin, it has been previously shown that mitochondrial dsRNA can activate these pathways64. When SGs are present, they may sequester these leaked nucleic acids, possibly explaining their frequent proximity to mitochondria, thereby dampening inflammatory signaling and perhaps even delaying SASP onset. 
Indeeed, several recent studies have shown that mitochondrial dsRNA can localize to SGs, either through increased mitochondrial permeability65 or stress-induced release through fibrosis66 or aging67. However, this raises an important issue: both fibrosis and aging are strongly linked to oxidative stress68,69, which in turn promotes mitochondrial dysfunction and nucleic acid leakage70. It therefore remains unclear whether the proposed mitochondria–SG–inflammation axis represents a general stress response, or whether it is specific to stresses that directly disrupt mitochondrial integrity.

\section{Materials and Methods}

\subsection{U-2 OS growth conditions}

Human osteosarcoma U2-OS cells were cultured in \gls{dmem} (high glucose) supplemented with 10\% fetal bovine serum. Cells were plated on glass coverslips in 24-well plates and used for all experimental procedures.

\subsection{Stress conditions}
To examine mRNA localization under stress, when 80\% confluency was reached, cells were exposed to 0.4 M sorbitol for 3 hours. To inhibit mitochondrial dsRNA leakage into the cytosol, 50 μM 4,4-diisothiocyanatostilbene-2,2-disulfonate (DIDS) (as previously described65) was applied either in combination with sorbitol or alone (control) for the same duration.

\subsection{Immunofluorescence}

Following stress treatments, cells were washed three times with PBS+/+. Fixation was performed with 4\% formaldehyde in PBS+/+ (pH 7.4) for 20 min at room temperature, after which cells were washed three additional times with PBS+/+. Permeabilization was carried out using 0.1\% Triton X-100 in PBS+/+ for 10 min at room temperature, followed by three washes in unsupplemented PBS.
Glass coverslips were then removed from the 24-well plates and transferred to a humid chamber prepared with 20 µL droplets of blocking buffer (5\% goat serum in PBS). Coverslips were inverted onto the droplets (cell side facing down) and incubated for 60 min at room temperature. Primary antibodies against HSP60 (mitochondria), J2 (dsRNA), and G3BP1 or TIA1 (stress granules) were diluted 1:200 in blocking buffer and applied to coverslips using the drop technique (10 µL per coverslip). Incubation proceeded for 2 h at room temperature in the humid chamber. Afterward, coverslips were washed three times in unsupplemented PBS.
Secondary antibodies were diluted 1:500 in blocking buffer together with DAPI (1:1000) and applied to coverslips (10 µL per coverslip) for 1 h in the dark at room temperature. Coverslips were then washed three times in unsupplemented PBS with agitation (5 min each wash). For mounting, 2 µL of Vectashield was placed on each microscope slide, and two coverslips were carefully placed cell-side down onto the drops. Coverslips were sealed with nail polish prior to imaging.

\subsection{Microscopy and image analysis}

For Figure 9 and Supplementary Figure 6, after fixation, cells were imaged on a ZEISS LSM 880 confocal microscope, using a 63x/1.4 Plan-Apochromat DIC oil immersion objective. For each condition, a z-stack of a single field of view (1024x1024px, with a px size of 0.13 µm) was acquired, comprising 10–20 Z-planes, with a 0.48 µm step size, according to Nyquist, with a scan zoom of 1.0. G3BP1 and TIA1 were excited using a 633 nm laser, HSP60 with a 561 nm laser, dsRNA J2 with a 488 nm laser, and DAPI with a 405 nm laser. Images were observed in FIJI, with maximum intensity thresholds defined for each channel using sorbitol-stressed samples. These same parameters were then applied to all other conditions to ensure consistency across images.
For Figure 10, the same samples from Figure 9 were re-acquired using a ZEISS LSM 980 with Airyscan 2 confocal microscope, with a 63x/1.4 Plan-Apochromat DIC oil immersion objective, using the Airyscan 2 in Super Resolution mode (Airyscan SR). Three fields of view were acquired per condition (2245x2245px, with a px size of 0.035 µm), with each z-stack comprising 35–65 Z-planes with a 0.130 µm step size, according to Nyquist, with a scan zoom of 1.7. G3BP1 was excited using a 639 nm laser, HSP60 with a 561 nm laser, dsRNA J2 with a 488 nm laser, and DAPI with a 405 nm laser.  Raw Airyscan data were processed using Airyscan Joint Deconvolution to improve resolution (Sample structure "Dense" and 5 Maximum Iterations). Basal autofluorescence levels were determined from a control sample incubated with secondary antibodies only and used as a minimal threshold for normalizing fluorescence signals in experimental images.
Colocalization analysis was performed in Imaris 10.2.0 CF, using the Coloc module, quantifying the percentage of overlapping voxels across different channels. Initially, colocalization between HSP60 (mitochondria) and J2 (dsRNA) was measured. Negative controls were used to define the minimum intensity threshold for colocalization of each channel. A new channel was subsequently generated, containing only the J2 signal that did not overlap with HSP60. This non-mitochondrial J2 signal was then analyzed for colocalization with G3BP1 (stress granules) under both unstressed and sorbitol-stressed conditions. All image processing and quantifications were performed on an image analysis and processing workstation with support from the Bioimaging Platform of the Gulbenkian Institute for Molecular Medicine Foundation. 
Software references are supplied in the Key Resources Table.

Regarding the statistical analysis of the data, colocalization percentages were initially compared between stressed and unstressed samples using Student’s t-tests. However, given the limited sample size, beta distributions were fitted to the three replicates per condition. From these fitted distributions, 1000 random samples were drawn and compared across several statistical metrics, including Cohen’s d, probability of difference (p-diff), Bhattacharyya coefficient (BC), and Area Under the Curve (AUC).


\section{Mitochondrial Preamibility is necessary for SG assembly}

To avoid confounding effects of oxidative stress, which is tightly linked to mitochondrial dysfunction and nucleic acid leakage70, we instead used hyperosmotic stress (through sorbitol exposure) in U2-OS cells, a widely employed model in SG research (reviewed in Figure 1B). Although hyperosmotic stress can impact mitochondrial function71, we reasoned that it would perturb mitochondria less directly than oxidative stress, allowing us to better assess stress granule–associated RNA dynamics. 
Previous work has shown that inhibition of either BAX/BAK or the mitochondrial porins VDAC1/2 prevents leakage of mitochondrial dsRNA into the cytoplasm65. We hypothesized that this mechanism would also apply during hyperosmotic stress and therefore included the VDAC1/2 inhibitor 4,4-diisothiocyanatostilbene-2,2-disulfonate  (DIDS)65 as an additional control. Detecting dsRNA within SGs under these conditions would strengthen the case that its origin is mitochondrial rather than from alternative sources. 
To test our model and directly visualize the putative SG–mitochondria–inflammation axis, we performed immunofluorescence using G3BP1 (SG marker, reviewed in Figure 1B), HSP60 (mitochondrial marker72), J2 (dsRNA marker65), and DAPI (nuclei).
Using the conditions described in T. Matheny 202150,51, we confirmed by 63× confocal imaging that U2-OS cells exposed to 0.4 M sorbitol for 3 h robustly formed SGs, visualized as G3BP1-circular foci, which were absent in unstressed controls (Figure 9). Strikingly, co-treatment with the VDAC inhibitor DIDS65 (50 μM) completely abrogated SG assembly (Figure 9). This finding suggests that mitochondrial permeability, or at least VDAC channel activity, is required for SG formation under hyperosmotic stress.
Notably, SGs have previously been reported to modulate mitochondrial permeability by regulating VDAC activity73. The observation that VDAC inhibition blocks SG assembly could point toward a reciprocal relationship, in which SGs both influence and depend on mitochondrial signaling.
At the same time, closer inspection showed that in both stressed and unstressed cells treated with DIDS, G3BP1 displayed a broadly similar diffuse distribution (Figure 9). This suggests that DIDS may also exert off-target or global cellular effects, making it difficult to attribute the phenotype exclusively to VDAC inhibition.
Nevertheless, the requirement for VDAC activity is consistent with a model in which mitochondrial-derived signals contribute to SG assembly. While leakage of mitochondrial nucleic acids remains a plausible mechanism, especially given VDAC’s reported role in their transport65, other damage-associated products (e.g., cytochrome, cardiolipin, metabolites) could also serve as triggers for inflammation/stress74 and sequential SG assembly. Alternatively, impaired import of mitochondrial factors might explain the effect.



\section{mtRNA is enriched in the SG}

Because SGs preferentially formed near mitochondria, as previously reported73, this imaging setup did not provide sufficient spatial resolution to directly quantify dsRNA colocalization within SGs.
To overcome the resolution limitations of conventional confocal microscopy, we re-imaged the same samples using Airyscan with joint deconvolution, which markedly improved spatial resolution. Because DIDS treatment altered the baseline distribution of G3BP1, we restricted our comparisons to stressed versus unstressed conditions only. Unexpectedly, at this resolution G3BP1 signal appeared as small, spherical foci in both stressed and unstressed cells, rather than the diffuse distribution typically observed in unstressed cells with conventional confocal imaging (Figure 10A). This made it more challenging to distinguish bona fide SGs from background G3BP1 signal. Nevertheless, given that these were the same samples as in our prior analysis, we proceeded with quantitative colocalization measurements.
First, we calculated the fraction of mitochondrial volume that colocalized with dsRNA volume (in voxels) (Figure 10A, in Teal). Unexpectedly, we observed higher mitochondrial–dsRNA overlap in unstressed cells compared to stressed conditions. This counterintuitive result suggests that at least in the case of hyperosmotic stress, dsRNA clearance may be enhanced through nucleases or autophagic pathways, or alter mitochondrial morphology in ways that reduce measured colocalization. An alternative explanation is that even under “unstressed” culture conditions, cells experience baseline stress that may not be sufficient to trigger SG formation but could still promote mitochondrial dsRNA leakage. Because SGs are thought to have protective functions7, their assembly during acute stress might reduce such leakage, independently of any direct uptake. Moreover, evidence from C. elegans suggests that inhibition of protein synthesis enhances mitochondrial function75. Since protein synthesis is typically suppressed under strong stress11, this could paradoxically increase mitochondrial activity and dsRNA production.
To accurately assess SG-associated uptake, we excluded dsRNAs overlapping both mitochondria and SGs before quantifying their colocalization with G3BP1. Using this filtered signal, we then measured the proportion of dsRNA volume colocalizing with G3BP1 signal (Figure 10A, in Magenta). Colocalization was higher under stress, both within mitochondria (Figure 10B), and in SGs, the latter being consistent with dsRNA sequestration into SGs (Figure 10C). To account for our limited sample size, we modeled the data with a beta distribution, and compared these distributions (v. Methods). Although statistical significance was not achieved due to low power, the analysis revealed a large effect size between stressed and unstressed conditions (Area under the curve (AUC) = 0.969, Cohen’s d = 2.847) (Figure 10C), and differences between the two distributions (Bhattacharyya coefficient (BC) = 0.247) supporting a robust difference.
Taken together, these observations provide strong support for our hypothesis that SGs sequester cytoplasmic dsRNA under hyperosmotic stress. However, given the technical limitations and modest sample size, these data should be interpreted as supportive rather than definitive, and additional experiments will be required to fully validate this model.


